\section{Data Processing / Cleaning}

The Python part of this exercise is available in \myhref{https://github.com/LosDelDGIIM/LosDelDGIIM.github.io/blob/main/subjects/Erasmus-DUE/GKI/Exercises/Exercise1.ipynb}{this Jupyter Notebook}.

\subsection*{Data Filtering and Transforming}

\begin{ejercicio}
    Imagine you have a list of all the students in a school. You only want to look at the students in the 5th grade (data filtering). Then, you want to calculate the average age of those 5th graders (data transforming). Explain why you need both steps.\\

    We first need to filter the data to focus only on the students in the 5th grade, as we are specifically interested in that group. This step allows us to narrow down our dataset to the relevant subset of students. After filtering, we need to transform the data by calculating the average age of those 5th graders. This transformation step is necessary to derive meaningful insights from the filtered data, as it provides us with a summary statistic (the average age) that can help us understand the characteristics of the 5th-grade students. Without filtering, we would be calculating the average age of all students, which would not give us the specific information we are looking for about the 5th graders.
\end{ejercicio}

\begin{ejercicio}
    You have a list of toys that were sold in a store, including the toy's name and the price it sold for.
    \begin{enumerate}
        \item Give an example of a question you could answer by only looking at toys that cost more than $20$ (data filtering).
        
        A question that could be answered by data filtering would be: ``Which toys that cost more than 20 euros were sold in the last month?'' (Here, we filter the data by price greater than 20 euros and potentially by the sales date).
        \item Give an example of something you could calculate or change about the toy prices (data transforming).
        
        An example of data transformation would be calculating the average selling price of all toys or converting all prices from euros to another currency (e.g., US dollars) based on a specific exchange rate. Another transformation could be creating a new column indicating whether a toy is classified as ``cheap'' (e.g., $<$ 15 euros), ``mid-priced'' (e.g., 15 euros $-$ 30 euros), or ``expensive'' ($>$ 30 euros) based on the selling price.
    \end{enumerate}
\end{ejercicio}


% // TODO: Corregir
\subsection*{Data Understanding}

\begin{ejercicio}
    What is the difference between quantitative (numerical) and qualitative (categorical) data?\\

    Quantitative data refers to numerical values that can be measured and quantified, such as height, weight, or age. Qualitative data, on the other hand, refers to categorical values that describe characteristics or attributes, such as color or gender. Quantitative data can be analyzed using mathematical operations, while qualitative data is typically analyzed using descriptive statistics or categorization.
\end{ejercicio}

\begin{ejercicio}
    What does the abbreviation \verb|NaN| mean in a dataset?

    \verb|NaN| stands for ``Not a Number'' and is used in datasets to represent missing or undefined values. It is also used to indicate that a value is not a valid number, such as the result of an undefined mathematical operation (e.g., division by zero). In data analysis, \verb|NaN| values need to be handled appropriately, either by filling them with a specific value or dropping them.
\end{ejercicio}

\begin{ejercicio}
    Briefly explain the difference between the mean and the median. Which value is more sensitive to outliers?

    The mean is the average of a set of numbers, calculated by summing all the values and dividing by the count of values. The median is the middle value in a sorted list of numbers. The mean is more sensitive to outliers because it can be significantly affected by extremely high or low values, while the median is less affected as it only considers the middle value(s) and not the magnitude of all values.
\end{ejercicio}

\begin{ejercicio}
    Calculate the quartiles for the following dataset: [10, 9, 1, 4, 3, 2, 5, 6, 6, 6, 3, 9].

    First, we need to sort the dataset:
    \begin{equation*}
    [1, 2, 3, 3, 4, 5, 6, 6, 6, 9, 9, 10]
    \end{equation*}

    The quartiles are then:
    \begin{align*}
        Q2 &= \dfrac{6 + 5}{2} = 5.5\\
        Q1 &= \dfrac{3 + 3}{2} = 3\\
        Q3 &= \dfrac{6 + 9}{2} = 7.5
    \end{align*}
\end{ejercicio}

\subsection*{Data Cleaning}

\begin{ejercicio}
    Name two reasons why data is often dirty in the real world.\\

    Data can be dirty due to:
    \begin{itemize}
        \item Human error during data entry, which can lead to typos, incorrect values, or missing data.
        \item Inconsistent data formats or structures, such as different date formats or varying units of measurement across datasets.
    \end{itemize}
\end{ejercicio}

\begin{ejercicio}
    What happens when you apply the \verb|dropna()| function to a DataFrame?

    The \verb|dropna()| function removes any rows (or columns, if specified) that contain missing values (\verb|NaN|) from the DataFrame. This is a common data cleaning step to ensure that the dataset does not contain incomplete records that could affect analysis or modeling.
\end{ejercicio}

\begin{ejercicio}
    Why do you scale (normalize) numerical data before feeding it into a machine learning model?\\

    Scaling (normalizing) numerical data is important because many machine learning algorithms are sensitive to the scale of the input features. If the features have different ranges, the algorithm may give more weight to those with larger values, leading to biased results. Scaling helps to ensure that all features contribute equally to the model's learning process and can improve the convergence speed and performance of the model.
\end{ejercicio}

\begin{ejercicio}
    Name two possible encodings and explain them using an example.\\

    They are explained in the page~\pageref{item:encoding}.
\end{ejercicio}